\documentclass[12pt, a4paper]{article}
\usepackage[utf8]{inputenc}
\usepackage[T1]{fontenc}
\usepackage[french]{babel}
\usepackage[margin=2.5cm]{geometry}
\usepackage{amsmath, amssymb, amsthm, mathtools}
\usepackage{listings}
\usepackage{xcolor}
\usepackage{tcolorbox}
\tcbuselibrary{theorems, skins, breakable}
\usepackage{booktabs}
\usepackage{array}
\usepackage{fancyhdr}
\usepackage{hyperref}
\usepackage{tikz}
\usetikzlibrary{shapes, positioning}

% ── Couleurs ──────────────────────────────────────────────────────────
\definecolor{suva_blue}{RGB}{30, 100, 180}
\definecolor{code_bg}{RGB}{245, 247, 250}
\definecolor{code_kw}{RGB}{0, 80, 200}
\definecolor{code_str}{RGB}{180, 40, 40}
\definecolor{code_cmt}{RGB}{80, 130, 80}
\definecolor{warn_bg}{RGB}{255, 248, 220}
\definecolor{success_bg}{RGB}{230, 255, 235}
\definecolor{dark_bg}{RGB}{25, 30, 50}

\lstdefinestyle{python}{
  language=Python,
  basicstyle=\ttfamily\small,
  keywordstyle=\color{code_kw}\bfseries,
  stringstyle=\color{code_str},
  commentstyle=\color{code_cmt}\itshape,
  backgroundcolor=\color{code_bg},
  frame=single, rulecolor=\color{gray!40},
  numbers=left, numberstyle=\tiny\color{gray},
  breaklines=true, tabsize=4, showstringspaces=false,
}
\lstdefinestyle{output}{
  basicstyle=\ttfamily\small,
  backgroundcolor=\color{gray!10},
  frame=single, rulecolor=\color{gray!50},
  breaklines=true,
}

\newtcolorbox{defbox}[2][]{
  colback=success_bg, colframe=green!50!black,
  fonttitle=\bfseries, title={Définition : #2}, #1
}
\newtcolorbox{theorembox}[2][]{
  colback=blue!5, colframe=suva_blue,
  fonttitle=\bfseries, title={Théorème : #2}, #1
}
\newtcolorbox{exercice}[2][]{
  colback=white, colframe=suva_blue!60,
  fonttitle=\bfseries, title={Exercice #2}, #1, breakable
}
\newtcolorbox{solution}[1][]{
  colback=gray!5, colframe=gray!40,
  fonttitle=\bfseries\itshape, title={Solution}, #1, breakable
}
\newtcolorbox{importantbox}[1][]{
  colback=suva_blue!10, colframe=suva_blue,
  fonttitle=\bfseries, title={Connexion avec le projet}, #1
}
\newtcolorbox{warningbox}[1][]{
  colback=warn_bg, colframe=orange!80!black,
  fonttitle=\bfseries, title={Attention}, #1
}

\newcommand{\Zq}{\mathbb{Z}_q}
\newcommand{\ZZ}{\mathbb{Z}}
\newcommand{\pmod}[1]{\ (\mathrm{mod}\ #1)}

\newcommand{\auteur}{Mohamed Lamine MBENGUE}
\newcommand{\specialites}{Sécurité · DevOps · Dev Full Stack Web \& Mobile}
\newcommand{\tel}{+221 78 178 44 36}
\newcommand{\email}{mbenguemohamedlamine2022@gmail.com}
\newcommand{\portfolio}{portfolio-alamine.web.app}

\pagestyle{fancy}
\fancyhf{}
\fancyhead[L]{\color{suva_blue}\textbf{SuVa --- Chapitre 1}}
\fancyhead[R]{\color{gray}Arithmétique Modulaire}
\fancyfoot[L]{\footnotesize\color{gray}\auteur}
\fancyfoot[C]{\thepage}
\fancyfoot[R]{\footnotesize\color{gray}\portfolio}
\renewcommand{\headrulewidth}{0.4pt}
\renewcommand{\footrulewidth}{0.3pt}

% ══════════════════════════════════════════════════════════════════════
\begin{document}

\begin{center}
  {\color{suva_blue}\rule{\linewidth}{2pt}}\\[0.5cm]
  {\Huge\bfseries\color{suva_blue} Chapitre 1}\\[0.3cm]
  {\large\bfseries Arithmétique Modulaire et Corps Finis}\\[0.2cm]
  {\small\color{gray} Fondement mathématique de toute la cryptographie du projet}\\[0.3cm]
  {\color{suva_blue}\rule{\linewidth}{2pt}}
\end{center}

\begin{tcolorbox}[colback=suva_blue!5, colframe=suva_blue!40,
  left=6pt, right=6pt, top=4pt, bottom=4pt]
\small
\begin{tabular}{ll}
  \textbf{Auteur} & \auteur \\
  \textbf{Spécialités} & \specialites \\
  \textbf{Tél} & \tel \quad \textbf{Email} : \email \\
  \textbf{Portfolio} & \href{https://\portfolio}{\texttt{\portfolio}} \\
\end{tabular}
\end{tcolorbox}

\tableofcontents
\newpage

%══════════════════════════════════════════════════════════════════════
\section{Arithmétique modulaire --- rappels}
%══════════════════════════════════════════════════════════════════════

\subsection{L'opération modulo}

\begin{defbox}{Modulo}
Pour deux entiers $a \in \ZZ$ et $m \in \ZZ_{>0}$, on dit que
$a \equiv r \pmod{m}$ si $m$ divise $a - r$.
L'entier $r$ est le \textbf{reste} de la division euclidienne de $a$ par $m$.

On note aussi $a \bmod m = r$ avec $r \in \{0, 1, \ldots, m-1\}$.
\end{defbox}

\textbf{Exemples~:}
\begin{align*}
  17 \bmod 5 &= 2  &\text{car } 17 = 3 \times 5 + 2\\
  -3 \bmod 7 &= 4  &\text{car } -3 = (-1) \times 7 + 4\\
  100 \bmod 13 &= 9 &\text{car } 100 = 7 \times 13 + 9
\end{align*}

\begin{warningbox}
En Python, \texttt{-3 \% 7 = 4} (correct). En C, le résultat peut être négatif.
Dans le projet, tous les calculs doivent donner un résultat dans $[0, q-1]$.
\end{warningbox}

\subsection{L'anneau $\mathbb{Z}_q$}

\begin{defbox}{L'anneau $\mathbb{Z}_q$}
$\mathbb{Z}_q = \{0, 1, 2, \ldots, q-1\}$ muni de l'addition et de la
multiplication modulo $q$. C'est un \textbf{corps} (corps fini) si et
seulement si $q$ est \textbf{premier}.
\end{defbox}

\begin{importantbox}
Dans le projet SuVa / Dilithium, le module est~:
$$q = 8\,380\,417$$
Ce nombre est premier (vérifie-le !). Tous les coefficients des polynômes
vivent dans $\mathbb{Z}_q = \{0, 1, \ldots, 8\,380\,416\}$.
\end{importantbox}

\subsection{Propriétés de l'addition et de la multiplication modulo $q$}

Les règles habituelles s'appliquent~:
\begin{align*}
  (a + b) \bmod q &\equiv ((a \bmod q) + (b \bmod q)) \bmod q\\
  (a \cdot b) \bmod q &\equiv ((a \bmod q) \cdot (b \bmod q)) \bmod q
\end{align*}

\textbf{Exemple avec $q = 8\,380\,417$~:}
\begin{lstlisting}[style=python]
q = 8_380_417

a = 5_000_000
b = 4_000_000

addition = (a + b) % q
print(f"({a} + {b}) mod q = {addition}")
# (5000000 + 4000000) mod 8380417 = 619583

multiplication = (a * b) % q
print(f"({a} * {b}) mod q = {multiplication}")
# Resultat dans [0, 8380416]
\end{lstlisting}

%══════════════════════════════════════════════════════════════════════
\section{L'inverse modulaire}
%══════════════════════════════════════════════════════════════════════

\subsection{Définition}

\begin{defbox}{Inverse modulaire}
L'inverse de $a$ modulo $q$ est l'entier $a^{-1} \in \{1, \ldots, q-1\}$
tel que~:
$$a \cdot a^{-1} \equiv 1 \pmod{q}$$
L'inverse existe si et seulement si $\gcd(a, q) = 1$.
Si $q$ est premier, tout $a \neq 0$ est inversible.
\end{defbox}

\subsection{Algorithme d'Euclide étendu}

\begin{theorembox}{Théorème de Bézout}
Pour tout $a, b$ tels que $\gcd(a, b) = d$, il existe $u, v \in \ZZ$
tels que~:
$$au + bv = d$$
Si $d = 1$, alors $u = a^{-1} \pmod{b}$.
\end{theorembox}

\begin{lstlisting}[style=python]
def xgcd(a, b):
    """
    Algorithme d'Euclide etendu.
    Retourne (gcd, u, v) tels que a*u + b*v = gcd.
    """
    if b == 0:
        return a, 1, 0
    g, u, v = xgcd(b, a % b)
    return g, v, u - (a // b) * v

def mod_inverse(a, q):
    """Inverse de a modulo q (q premier)."""
    g, u, v = xgcd(a % q, q)
    assert g == 1, f"{a} n'est pas inversible modulo {q}"
    return u % q

# Test
q = 8_380_417
a = 42
inv = mod_inverse(a, q)
print(f"42^(-1) mod q = {inv}")
print(f"Verification : 42 * {inv} mod q = {(42 * inv) % q}")  # Doit etre 1
\end{lstlisting}

\subsection{Méthode rapide : petit théorème de Fermat}

\begin{theorembox}{Petit théorème de Fermat}
Si $q$ est premier et $\gcd(a, q) = 1$, alors~:
$$a^{q-1} \equiv 1 \pmod{q}$$
Donc~: $a^{-1} \equiv a^{q-2} \pmod{q}$
\end{theorembox}

\begin{lstlisting}[style=python]
def mod_inverse_fermat(a, q):
    """Inverse via le petit theoreme de Fermat (q premier)."""
    return pow(a, q - 2, q)  # pow(a, e, m) = a^e mod m (Python integre)

q = 8_380_417
a = 42
inv = mod_inverse_fermat(a, q)
print(f"42^(-1) mod q = {inv}")
print(f"Verification  : {(42 * inv) % q}")  # 1
\end{lstlisting}

\begin{importantbox}
La fonction \texttt{pow(a, q-2, q)} de Python est \textbf{très efficace}~:
elle utilise l'exponentiation rapide en $O(\log q)$ multiplications.
C'est ce qui est utilisé dans le projet pour les inversions modulaires.
\end{importantbox}

%══════════════════════════════════════════════════════════════════════
\section{Réduction symétrique (centred reduction)}
%══════════════════════════════════════════════════════════════════════

\subsection{Motivation}

Dans Dilithium, les coefficients ne vivent pas toujours dans $[0, q-1]$.
Parfois on préfère les représenter dans $[-(q-1)/2, (q-1)/2]$.

\begin{defbox}{Réduction symétrique}
La réduction symétrique de $a$ modulo $q$ est~:
$$a_{\mathrm{sym}} = \begin{cases}
  a \bmod q & \text{si } (a \bmod q) \leq (q-1)/2\\
  (a \bmod q) - q & \text{sinon}
\end{cases}$$
Résultat dans $[-(q-1)/2, (q-1)/2]$.
\end{defbox}

\begin{lstlisting}[style=python]
# Fichier : utilities/utils.py dans le projet
def reduce_mod_pm(a, q):
    """Reduction symmetrique de a modulo q."""
    r = a % q
    if r > (q - 1) // 2:
        r -= q
    return r

q = 8_380_417
# Exemples
print(reduce_mod_pm(8_380_416, q))  # = -1 (pas q-1)
print(reduce_mod_pm(1, q))          # =  1
print(reduce_mod_pm(q // 2 + 1, q)) # = -(q//2)
\end{lstlisting}

\begin{importantbox}
Cette fonction apparaît dans \texttt{utilities/utils.py}.
Elle est utilisée dans la vérification de normes ($\|z\|_\infty < \gamma_1 - \beta$).
\end{importantbox}

%══════════════════════════════════════════════════════════════════════
\section{Décomposition power-of-2 (HighBits / LowBits)}
%══════════════════════════════════════════════════════════════════════

C'est un mécanisme crucial dans Dilithium pour compresser la clé publique.

\subsection{Principe}

\begin{defbox}{Décomposition modulaire}
Pour $r \in \ZZ_q$ et $\alpha$ pair divisant $q - 1$, on écrit~:
$$r = r_1 \cdot \alpha + r_0$$
avec $r_0 \in (- \alpha/2,\ \alpha/2]$ (partie basse, \emph{LowBits})
et $r_1 = (r - r_0) / \alpha$ (partie haute, \emph{HighBits}).
\end{defbox}

\textbf{Pourquoi~?} Dans Dilithium, $t = A s_1 + s_2$ est un vecteur de polynômes.
On décompose $t = t_1 \cdot 2^d + t_0$. La clé publique ne publie que $t_1$
(bits hauts), économisant $d = 13$ bits par coefficient. $t_0$ reste secret.

\begin{lstlisting}[style=python]
# Extrait de utilities/utils.py
def decompose(r, alpha, q=8_380_417):
    """
    Decomposes r = r1 * alpha + r0
    r0 in (-alpha/2, alpha/2]
    """
    r = r % q
    r0 = r % alpha
    if r0 > alpha // 2:
        r0 -= alpha
    r1 = (r - r0) // alpha
    if r1 == (q - 1) // alpha + 1:
        r1 = 0
        r0 -= 1
    return r1, r0

def high_bits(r, alpha, q=8_380_417):
    r1, _ = decompose(r, alpha, q)
    return r1

def low_bits(r, alpha, q=8_380_417):
    _, r0 = decompose(r, alpha, q)
    return r0

# Test
q = 8_380_417
d = 13
alpha = 1 << d  # 2^13 = 8192

r = 100_000
r1, r0 = decompose(r, alpha, q)
print(f"r  = {r}")
print(f"r1 = {r1}  (partie haute, dans la cle publique)")
print(f"r0 = {r0}  (partie basse, secrete)")
print(f"Verification : r1 * {alpha} + r0 = {r1 * alpha + r0}")
# Doit etre congruent a r mod q
\end{lstlisting}

%══════════════════════════════════════════════════════════════════════
\section{Le vecteur hint --- make\_hint et use\_hint}
%══════════════════════════════════════════════════════════════════════

\subsection{Motivation}

Lors de la vérification, on recalcule $w' = Az - ct_1 \cdot 2^d$ et on
veut retrouver $w_1 = \mathrm{HighBits}(Ay)$ d'origine.
Le problème~: $Az - ct_1 \cdot 2^d \approx Ay$ mais pas exactement.
Le \textbf{hint} $h$ encode la correction nécessaire de manière compacte.

\begin{defbox}{Hint}
Pour $r \in \ZZ_q$ et $r + z \in \ZZ_q$ avec $z$ petit~:
$$h = \mathrm{MakeHint}(z, r, \alpha) \in \{0, 1\}$$
$$r_1 = \mathrm{UseHint}(h, r + z, \alpha)$$
de sorte que $r_1 = \mathrm{HighBits}(r, \alpha)$.
\end{defbox}

\begin{lstlisting}[style=python]
# Extrait de utilities/utils.py
def make_hint(z, r, alpha, q=8_380_417):
    """
    Cree un hint indiquant si HighBits(r, alpha) != HighBits(r+z, alpha).
    Retourne 0 ou 1.
    """
    r1 = high_bits(r, alpha, q)
    v1 = high_bits(r + z, alpha, q)
    return int(r1 != v1)

def use_hint(h, r, alpha, q=8_380_417):
    """
    Utilise le hint pour recuperer HighBits(r - z, alpha)
    sans connaitre z.
    """
    m = (q - 1) // alpha
    r1, r0 = decompose(r, alpha, q)
    if h == 1 and r0 > 0:
        return (r1 + 1) % m
    if h == 1 and r0 <= 0:
        return (r1 - 1) % m
    return r1

# Demonstration
q = 8_380_417
alpha = 2 * ((q - 1) // 88)  # = 2 * gamma_2 dans Dilithium

r = 500_000   # valeur originale
z = 90_000    # petite perturbation (typiquement cs2 + ct0)

r1_original = high_bits(r, alpha, q)
h = make_hint(z, r, alpha, q)
r1_recovered = use_hint(h, r + z, alpha, q)

print(f"r1 original  = {r1_original}")
print(f"hint h       = {h}")
print(f"r1 recupere  = {r1_recovered}")
print(f"Correct ?    = {r1_original == r1_recovered}")
\end{lstlisting}

%══════════════════════════════════════════════════════════════════════
\section{Vérification de la norme infinie}
%══════════════════════════════════════════════════════════════════════

\begin{defbox}{Norme infinie}
Pour un polynôme $f = \sum_{i=0}^{n-1} f_i \cdot x^i \in R_q$, la norme
infinie (après réduction symétrique) est~:
$$\|f\|_\infty = \max_{i=0}^{n-1} |f_i|$$
\end{defbox}

\begin{importantbox}
Dans Dilithium, les rejections sampling vérifient~:
\begin{itemize}
  \item $\|z\|_\infty < \gamma_1 - \beta$ \hfill ($\gamma_1 = 2^{17}$, $\beta = 78$)
  \item $\|w_0 - cs_2\|_\infty < \gamma_2 - \beta$ \hfill ($\gamma_2 = (q-1)/88$)
  \item $\|ct_0\|_\infty < \gamma_2$
\end{itemize}
Si une de ces conditions échoue, le signe recommence.
\end{importantbox}

\begin{lstlisting}[style=python]
def check_norm_bound(poly_coeffs, bound, q=8_380_417):
    """
    Verifie si |coefficient| >= bound pour au moins un coefficient.
    Retourne True si la norme depasse la borne (==> rejeter).
    """
    for coeff in poly_coeffs:
        # Reduction symmetrique
        c = coeff % q
        if c > (q - 1) // 2:
            c -= q
        if abs(c) >= bound:
            return True
    return False

# Test
import random
q = 8_380_417
gamma_1 = 1 << 17    # 131072
beta = 78
bound = gamma_1 - beta  # 130994

# Un vecteur z valide (petits coefficients)
z_valid = [random.randint(-bound + 1, bound - 1) for _ in range(256)]
print(f"z_valid depasse borne ? {check_norm_bound(z_valid, bound)}")  # False

# Un vecteur z invalide (trop grand)
z_bad = z_valid.copy()
z_bad[42] = gamma_1  # coefficient trop grand
print(f"z_bad depasse borne  ? {check_norm_bound(z_bad, bound)}")   # True
\end{lstlisting}

%══════════════════════════════════════════════════════════════════════
\section{Exercices}
%══════════════════════════════════════════════════════════════════════

\begin{exercice}{1 --- Calculs modulo q}
Soit $q = 8\,380\,417$.

\begin{enumerate}[label=\alph*)]
  \item Calcule $(7\,000\,000 + 2\,000\,000) \bmod q$.
  \item Calcule $(-1) \bmod q$.
  \item Calcule $q^2 \bmod q$.
  \item Trouve $a^{-1} \bmod q$ pour $a = 256$.
  \item Vérifie que $256 \times a^{-1} \equiv 1 \pmod{q}$.
\end{enumerate}
\end{exercice}

\begin{solution}
\begin{lstlisting}[style=python]
q = 8_380_417

# a)
print((7_000_000 + 2_000_000) % q)   # 619583

# b)
print((-1) % q)                        # 8380416

# c)
print((q**2) % q)                      # 0  (q^2 = q * q)

# d)
inv_256 = pow(256, q - 2, q)
print(f"256^(-1) mod q = {inv_256}")

# e)
print((256 * inv_256) % q)             # 1
\end{lstlisting}
\end{solution}

\begin{exercice}{2 --- Reduction symétrique}
\begin{enumerate}[label=\alph*)]
  \item Calcule la réduction symétrique de $8\,380\,416$ modulo $q$.
  \item Calcule la réduction symétrique de $4\,000\,000$ modulo $q$.
  \item Quelle est la plage de la réduction symétrique pour $q = 8\,380\,417$ ?
\end{enumerate}
\end{exercice}

\begin{solution}
\begin{lstlisting}[style=python]
q = 8_380_417

def sym(a, q):
    r = a % q
    return r - q if r > (q-1)//2 else r

print(sym(8_380_416, q))  # -1
print(sym(4_000_000, q))  # 4000000 (deja dans la plage)
# Plage : [-(q-1)//2, (q-1)//2] = [-4190208, 4190208]
\end{lstlisting}
\end{solution}

\begin{exercice}{3 --- Décomposition et hint (le plus important)}
Soit $q = 8\,380\,417$, $\alpha = 2\gamma_2$ avec $\gamma_2 = (q-1)/88$.

\begin{enumerate}[label=\alph*)]
  \item Calcule $\alpha$.
  \item Décompose $r = 3\,000\,000$ : calcule $r_1, r_0$.
  \item Vérifie que $r_1 \times \alpha + r_0 \equiv r \pmod{q}$.
  \item Soit $z = 50\,000$. Calcule $h = \mathrm{MakeHint}(z, r, \alpha)$.
  \item Vérifie que $\mathrm{UseHint}(h, r+z, \alpha) = \mathrm{HighBits}(r, \alpha)$.
\end{enumerate}
\end{exercice}

\begin{solution}
\begin{lstlisting}[style=python]
q = 8_380_417
gamma_2 = (q - 1) // 88
alpha = 2 * gamma_2
print(f"alpha = {alpha}")  # 190464

# b) Decomposer r = 3_000_000
r = 3_000_000
r0 = r % alpha
if r0 > alpha // 2:
    r0 -= alpha
r1 = (r - r0) // alpha
print(f"r1 = {r1}, r0 = {r0}")

# c) Verification
print(f"r1*alpha + r0 = {r1 * alpha + r0}")  # == r

# d) Hint
z = 50_000
h = make_hint(z, r, alpha, q)
print(f"hint h = {h}")

# e) UseHint
r1_rec = use_hint(h, r + z, alpha, q)
r1_orig = high_bits(r, alpha, q)
print(f"r1_orig = {r1_orig}, r1_rec = {r1_rec}")
print(f"Correct : {r1_orig == r1_rec}")
\end{lstlisting}
\end{solution}

\begin{exercice}{4 --- Exploration du code source}
Ouvre le fichier \texttt{utilities/utils.py} du projet.
\begin{enumerate}[label=\alph*)]
  \item Identifie la différence entre \texttt{make\_hint} et \texttt{make\_hint\_optimised}.
        (Indice~: Dilithium 5.1 vs spec NIST)
  \item Quel est l'avantage de \texttt{make\_hint\_optimised} ?
  \item À quel endroit dans \texttt{Falcon\_dilithium.py} est appelé \texttt{make\_hint\_optimised} ?
\end{enumerate}
\end{exercice}

\end{document}
